\subsection{复测度}

定义 6. --- 把复向量空间 \(\mathcal{K}_{\mathbf{C}}(\mathrm{T})\) 上的每个连续线性形式都称为 T 上的复测度。 \(^{2}\)

于是，T 上复测度所成的空间 \(\mathcal{M}_{\mathbf{C}}(\mathrm{T})\) 就是 Hausdorff 局部凸空间 \(\mathcal{K}_{\mathbf{C}}(\mathrm{T})\) 的对偶。

若 \(m\) 是 \(\mathrm{T}\) 上的复测度，它在 \(\mathcal{H}(\mathrm{T})\) 上的限制是 \(\mathrm{T}\) 上取值于 \(\mathbf{C}\) 的向量测度（看作 \(\mathbf{R}\) 上的向量空间）；

m 由这一限制确定，因为若 \(f = f_{1} + if_{2} \in \mathcal{K}_{\mathbf{C}}(\mathrm{T})\) ，则 f 的实部 \(f_{1}\) 与虚部 \(f_{2}\) 都属于 \(\mathcal{K}(\mathrm{T})\) ，且 \(m(f) = m(f_{1}) + im(f_{2})\) 。反过来，对 T 上取值于 C 的每个向量测度 \(m_{0}\) ，公式 \(m(f) = m_{0}(f_{1}) + im_{0}(f_{2})\) 定义一个复测度 m，它是 T 上唯一的在 \(\mathcal{K}(\mathrm{T})\) 上的限制等于 \(m_{0}\) 的测度。因此，今后我们把复测度与它在 \(\mathcal{K}(\mathrm{T})\) 上的限制等同起来；这样的测度 m 具有形式 \(m = \mu_{1} + i\mu_{2}\) ，其中 \(\mu_{1}\) 与 \(\mu_{2}\) 是 T 上的两个实测度，分别称为 m 的实部与虚部。m 的支集是 \(\mu_{1}\) 与 \(\mu_{2}\) 的支集的并。已知 m 是可优超的（No.~3, 命题 4 的推论）；我们把正测度 \(|m|\) 称为 m 的绝对值，它对应于 C 上的绝对值 \(|x_{1} + ix_{2}| = \sqrt{x_{1}^{2} + x_{2}^{2}}\) 。我们有 \(|m| = (\mu_{1}^{2} + \mu_{2}^{2})^{1/2}\) （No.~4, 命题 9 之后的注记）, \(^{3}\) 以及 \(|\mu_{1}| \leqslant |m|\) ， \(|\mu_{2}| \leqslant |m|\) ， \(|m| \leqslant |\mu_{1}| + |\mu_{2}|\) ；此外，m 是以 \(|m|\) 为基的测度，且可写 \(m = h \cdot |m|\) ，其中 \(h \in \mathcal{L}_{\mathbf{C}}^{\infty}(|m|)\) ，且 \(|h(t)| = 1\) 关于 \(|m|\) 局部几乎处处成立（No.~4, 命题 9）。 \(^{4}\) m 与 \(|m|\) 的支集相同。

对每个映射 \(\mathbf{f}\) （由 \(\mathrm{T}\) 到复 Banach 空间 \(\mathbf{F}\) 中的映射，关于 \(|m|\) 本质可积），可以定义（No.~7） \(\mathbf{f}\) 关于 \(m\) 的积分（对应于 \(\mathbf{R}\) -双线性映射 \((\mathbf{x},\lambda)\mapsto\lambda\mathbf{x}\) ，由 \(\mathbf{F}\times\mathbf{C}\) 到 \(\mathbf{F}\) 中），把它记作 \(\int\mathbf{f}dm\) ；由命题 11 的唯一性性质立即得出（采用前面的记号） \(\int\mathbf{f}dm=\int\mathbf{f}d\mu_{1}+i\int\mathbf{f}d\mu_{2}=\int\mathbf{f}hd|m|\) 。因此有 \(m(f)=\int fdm\) （对 \(f\in\mathcal{H}_{\mathbf{C}}(\mathrm{T})\) ）。我们说 \(\mathbf{f}\) 关于 \(m\) 本质可积，意思就是它关于 \(|m|;^{5}\) 本质可积； \(m\) -可积、 \(m\) -可测、局部 \(m\) -可积、 \(m\) -可忽略或局部 \(m\) -可忽略的映射都类似地定义。我们写

\[
\mathcal {L} _ {\mathrm{F}} ^ {p} (\mathrm{T}, m) \left(\mathrm{resp.} \overline {{\mathcal {L}}} _ {\mathrm{F}} ^ {p} (\mathrm{T}, m), \mathrm{L} _ {\mathrm{F}} ^ {p} (T, m)\right)
\]

代替 \(\mathcal{L}_{\mathrm{F}}^{p}(\mathrm{T},|m|)\) （相应地， \(\overline{\mathcal{L}}_{\mathrm{F}}^{p}(\mathrm{T},|m|)\) ， \(\mathrm{L}_{\mathrm{F}}^{p}(\mathrm{T},|m|)\) ）；它们都是复向量空间。

为使 f 为 m-可积（相应地，本质 m-可积），必须且只需 f 关于四个测度 \(\mu_{1}^{+}, \mu_{1}^{-}, \mu_{2}^{+}, \mu_{2}^{-}\) 中的每一个都可积（相应地，本质可积），这要依据 \(|m|, |\mu_{1}|, |\mu_{2}|\) 之间的诸不等式以及关系式 \(|\mu_{k}| = \mu_{k}^{+} + \mu_{k}^{-}\) （Ch. V, §2, No.~2, 命题 3 及其推论 1）。

若 f 本质 m-可积（相应地，m-可积），则 \(|f|\) 本质 \(|m|\) -可积（相应地， \(|m|\) -可积），且由命题 11 得出

\[
\left| \int \mathbf {f} d m \right| \leqslant \int | \mathbf {f} | d | m |.\tag{7}
\]

设 F 与 G 为两个复 Banach 空间，u 为 F 到 G 中的连续线性映射。若 f 是 T 到 F 中的本质 m-可积（相应地，m-可积）映射，则 \(u \circ f\) 本质 m-可积（相应地，m-可积），且 \(\int(u \circ f) dm = u (\int f dm)\) ；这由前文以及关于本质 \(|m|\) -可积函数的类似命题立即得出（Ch. IV, §4, No.~2, 定理 1 以及 Ch. V, §1, No.~3, 引理与定义 3）。

设 \(m\) 是 \(T\) 上的复测度， \(h\) 是局部 \(m\) -可积的复函数。对每个函数 \(f \in \mathcal{K}_{\mathbf{C}}(T)\) ，函数 \(fh\) 是 \(m\) -可积的，且映射 \(f \mapsto \int fh dm\) 是一个复测度，记作 \(h \cdot m\) ，称为以 \(h\) 为密度（关于 \(m\) ）的测度。若 \(m = g \cdot |m|\) ，则显然 \(h \cdot m = hg \cdot |m|\) ；此外，由于 \(|g(t)| = 1\) 关于 \(|m|\) 局部几乎处处成立，为使 \(\mathbf{f}\) 关于 \(n = h \cdot m\) 本质可积，必须且只需 \(\mathbf{f}h\) 本质 \(m\) -可积，且这时有 \(\int \mathbf{f} dn = \int (\mathbf{f}h) dm\) 。此外还有 \(|h \cdot m| = |h| \cdot |m|\) 。我们也把每个形如 \(h \cdot m\) 的测度称为以 \(m\) 为基的复测度；若两个复测度 \(m, m'\) 中的每一个关于另一个都具有密度，亦即 \(m' = h \cdot m\) （其中 \(h\) 局部 \(m\) -可积，且 \(h(t) \neq 0\) 关于 \(|m|\) 局部几乎处处成立），则称这两个复测度等价。显然 \(m\) 与 \(|m|\) 等价，且为使 \(m\) 与 \(m'\) 等价，必须且只需 \(|m|\) 与 \(|m'|\) 等价。

若 m 与 \(m'\) 是 T 上的两个复测度，f 是取值于复 Banach 空间 F 的函数，关于 m 与 \(m'\) 都本质可积（相应地，可积），则对任意复数 \(\lambda\) 与 \(\lambda'\) ，f 关于 \(\lambda m + \lambda'm'\) 本质可积（相应地，可积），且

\[
\int \mathbf {f} d (\lambda m + \lambda^ {\prime} m ^ {\prime}) = \lambda \int \mathbf {f} d m + \lambda^ {\prime} \int \mathbf {f} d m ^ {\prime}.
\]

因为这由第 7 目命题 11 的推论得出。

此外，由定义得出

\[
\left| \lambda m + \lambda^ {\prime} m ^ {\prime} \right| \leqslant \left| \lambda \right| \cdot \left| m \right| + \left| \lambda^ {\prime} \right| \cdot \left| m ^ {\prime} \right|.
\]

我们把复测度 m 的共轭测度称为复测度 \(\overline{m}\) ，它由 \(\overline{m}(f) = \overline{m(\overline{f})}\) （对 \(f \in \mathcal{K}_{\mathbf{C}}(\mathrm{T})\) ）定义。若 \(m = \mu_{1} + i\mu_{2}\) （其中 \(\mu_{1}\) 与 \(\mu_{2}\) 是实测度），则 \(\overline{m} = \mu_{1} - i\mu_{2}\) 且 \(|\overline{m}| = |m|\) ；若 \(m = h \cdot |m|\) ，则 \(\overline{m} = \overline{h} \cdot |m|\) 。若 f 是取复值的本质 m-可积（相应地， \(m\) -可积）函数，则 \(\overline{f}\) 本质 \(\overline{m}\) -可积（相应地， \(\overline{m}\) -可积），且

\[
\int \overline {{f}} d \overline {{m}} = \overline {{\int f d m}}.
\]

命题 12. --- 设 m 是 T 上的复测度，p 与 q 是共轭指数（Ch. IV, §6, No.~4）。双线性型 \((f,g)\mapsto\int fg dm\) 在乘积空间 \(\mathcal{L}_{\mathbf{C}}^{p}(m)\times\mathcal{L}_{\mathbf{C}}^{q}(m)\) 上有定义且连续；不等式 \(\left|\int fg dm\right|\leqslant\mathrm{N}_{p}(f)\mathrm{N}_{q}(g)\) 成立，且 \(\mathrm{N}_{q}(g)\) 是 \(\mathrm{L}_{\mathbf{C}}^{p}(m)\) 上由线性形式 \(f\mapsto\int fg dm\) 通过过渡到商得到的连续线性形式的范数。

此外，若 \(1 \leqslant p < +\infty\) ，则复向量空间 \(\mathcal{L}_{\mathbf{C}}^{p}(m)\) 上的每个连续线性形式都具有 \(f \mapsto \int fg dm\) 的类型，其中 \(g\) 是 \(\mathcal{L}_{\mathbf{C}}^{q}(m)\) 中的一个函数，它在 \(\mathrm{L}_{\mathbf{C}}^{q}(m)\) 中的类被唯一确定。

由于 \(m = h \cdot |m|\) （其中 \(|h(t)| = 1\) 局部几乎处处成立），第一个断言由 Hölder 不等式（Ch. IV, §6, No.~4, 定理 2）立即得出；第二个断言来自 Ch. IV, §6, No.~4 的命题 3。最后，若 \(u\) 是 \(\mathcal{L}_{\mathbf{C}}^{p}\) 上的连续线性形式，则它到子空间 \(\mathcal{L}^p\) （ \(\mathcal{L}_{\mathbf{C}}^{p}\) 的（实）子空间）上的限制是 \(\mathcal{L}^p\) 到 \(\mathbf{C}\) 中的连续 R-线性映射；于是若 \(1 \leqslant p < +\infty\) ，它便具有 \(f \mapsto \int fg_1 d|m| + i \int fg_2 d|m|\) 的类型，其中 \(g_1\) 与 \(g_2\) 属于 \(\mathcal{L}^q\) （Ch. V, §5, No.~8, 定理 4）；令 \(g = (g_1 + ig_2)h^{-1}\) ，即得最后一个断言。

